% Section 3: Specification I — The Definition of "AIF"

\section{Specification I: The Definition of ``AIF''}
\label{sec:aif}

\subsection{The Legal Text}
\label{sec:aif-text}

AIFMD Article 4(1)(a) provides:

\begin{quote}
``\,`AIFs' means collective investment undertakings, including investment compartments thereof, which:
\begin{enumerate}
\renewcommand{\labelenumi}{(\roman{enumi})}
  \item raise capital from a number of investors, with a view to investing it in accordance with a defined investment policy for the benefit of those investors; and
  \item do not require authorisation pursuant to Article 5 of Directive 2009/65/EC [the UCITS Directive].''
\end{enumerate}
\end{quote}

\noindent This deceptively simple definition determines whether an investment vehicle falls within the AIFMD regulatory perimeter.

\subsection{Interpretive Challenges}
\label{sec:aif-interpretation}

Before encoding, we must resolve ambiguities:

\begin{enumerate}
  \item \textbf{``A number of investors''}: Does this mean $>1$? ESMA guidance suggests context-dependent analysis, considering whether the vehicle was designed to pool capital. We encode as \texttt{number\_of\_investors >= 2} but flag this as an interpretive choice.
  
  \item \textbf{``Defined investment policy''}: What counts as ``defined''? A single-asset SPV has a very narrow policy. We encode as a boolean input, acknowledging that determining this value requires human judgment.
  
  \item \textbf{``For the benefit of those investors''}: This is tautological for well-formed funds---no fund is structured to harm its investors. We encode as a boolean input defaulting to true.
  
  \item \textbf{Exclusions}: Article 2 excludes holding companies, securitisation SPVs, pension funds, etc. These require separate scope checks.
\end{enumerate}

\subsection{The Catala Specification}
\label{sec:aif-spec}

We define a structure capturing fund characteristics:

\begin{lstlisting}
declaration structure FundCharacteristics:
  data is_collective_investment_undertaking content boolean
  data raises_capital_from_investors content boolean
  data number_of_investors content integer
  data has_defined_investment_policy content boolean
  data invests_for_investor_benefit content boolean
  data requires_ucits_authorization content boolean
\end{lstlisting}

The core AIF definition scope:

\begin{lstlisting}
declaration scope AIFDefinition:
  input fund content FundCharacteristics
  output qualifies_as_aif content boolean
  output classification_reasoning content text

scope AIFDefinition:
  definition qualifies_as_aif equals
    fund.is_collective_investment_undertaking
    and fund.raises_capital_from_investors
    and fund.number_of_investors >= 2
    and fund.has_defined_investment_policy
    and fund.invests_for_investor_benefit
    and (not fund.requires_ucits_authorization)
\end{lstlisting}

\subsection{Article 2 Exclusions}
\label{sec:aif-exclusions}

Article 2(3) excludes various entities. We implement these as exceptions:

\begin{lstlisting}
scope AIFExclusions:
  # Default: apply base status
  definition final_aif_status equals base_aif_status
  definition exclusion_applied equals "No exclusion applied"

  # Article 2(3)(a) - Holding companies
  exception definition final_aif_status equals false
    under condition
      exclusions.exclusion_category = Some HoldingCompany
      and exclusions.is_holding_company_per_article_4
    consequence definition exclusion_applied equals
      "Excluded as holding company under Article 2(3)(a)"
\end{lstlisting}

The exception mechanism explicitly encodes the legal hierarchy: the base rule applies unless an exclusion condition is met.

\subsection{Interpretive Choices Documented}
\label{sec:aif-choices}

\begin{table*}[t]
\centering
\begin{tabular}{@{}llll@{}}
\toprule
Element & Choice Made & Alternative & Rationale \\
\midrule
``Number of investors'' & $\geq 2$ & Context-dependent & Literal reading \\
``Defined investment policy'' & Boolean input & Structured criteria & Human judgment \\
``For the benefit of'' & Boolean input & Explicit test & Tautological \\
\bottomrule
\end{tabular}
\caption{Interpretive choices in AIF definition encoding}
\label{tab:aif-choices}
\end{table*}

\subsection{What the Specification Captures and Does Not Capture}
\label{sec:aif-limits}

\paragraph{Captured:}
\begin{itemize}
  \item The logical structure of the AIF definition
  \item The hierarchy: base definition $\rightarrow$ exclusions
  \item Traceability: which rule determined the outcome
\end{itemize}

\paragraph{Not Captured:}
\begin{itemize}
  \item The judgment required to determine ``defined investment policy''
  \item Context-dependent interpretation of ``number of investors''
  \item Interaction with national transposition variations
  \item Temporal aspects (when was capital raised? when is status assessed?)
\end{itemize}

The specification surfaces \textit{where} human judgment is required. It does not replace that judgment.
